\section{Complete source banks and paid recovery}\label{sec:source-engine}

This section proves the common finite statement used in the later exclusions.
The bank must cover every surviving raw source. A list of selected return
endpoints, without the state and unselected work dependencies, does not suffice.

\subsection{Fixed geometries and complete return families}

Fix integers $d,b\ge1$, and put $L=d^2b$, $W=db$ and
$\Omega=\{0,\ldots,L-1\}$. The finite results allow arbitrary $b$; the
application to~(T) takes $b=b(d)=\max(1,\lceil\log_2d\rceil)$. A
\emph{fixed complete geometry} consists of the input-head and work-head
walks on $\{0,\ldots,T-1\}$, the halting time, and the distinct increasing
append times $t_0<\cdots<t_{L-1}$, where $T\le cL$. The walks move by at
most one cell per step. For the Lean model, extend each walk constantly
after $T-1$; this preserves its unit-step property and all hops before $T$.
The output at time $t_r$ demands input coordinate
\begin{equation}\label{eq:transpose-demand}
 \pi(jW+ib+k)=iW+jb+k
 \qquad(0\le i,j<d,\ 0\le k<b).
\end{equation}
The geometry can be an oblivious schedule or one trace shared by a set
of inputs. We make no assertion that an arbitrary adaptive transposer has
a trace containing nearly all its input information.

The dependency graph has the state edges $t\to t-1$ and every work-head
previous-visit edge $u\to p$, where $p$ is the last earlier visit of that
head to its cell at time~$u$. A work edge of length one coincides with a
state edge. Coincident edges may be identified for reachability, but an
edge is retained whenever any physical dependency supplies it. Input
reads are sources, not previous-visit edges. The complete value at a time
contains its state after the step, its written work symbols and any output
bit. Each such value has at most $\beta$ bits of capacity, for a fixed
$\beta>0$; this includes all data needed by later steps.

A \emph{genie} is a set $\Gamma$ of time vertices and its transcript
$G_\Gamma$ consists of those complete values. Thus
$H(G_\Gamma)\le\beta\card\Gamma$. A backward path
$v_0\to\cdots\to v_s$ survives the genie if
$v_1,\ldots,v_s\notin\Gamma$. In particular, its root is retained even
when $v_0\in\Gamma$: we retain its direct read and every incoming
dependency whose lower endpoint is unpaid. The full source relation at
an output contains every input coordinate read at the end of such a path,
including a path of length zero. The usual recursive evaluation of the
dependency graph shows that this full source relation and $G_\Gamma$
determine the output. This is the cone convention of
\cite[Section~4.1]{FourthPaper2026}.

Let $\mathcal C\subseteq\{0,1\}^{\Omega}$ be a nonempty set correctly
transposed on the fixed geometry, and write
\begin{equation}\label{eq:correctness-deficit}
 \Delta=L-\log_2\card{\mathcal C}.
\end{equation}
All selections of families, time maps and paid vertices below depend on
the fixed geometry, not on the input values within $\mathcal C$. We use
ordinary base-two Shannon entropy only for finite random variables.

Write $x(t)$ for the input-head position clamped to $[0,L-1]$. This does
not increase its total variation. A read outside the input is a known
blank; including the bit at the clamped position supplies extra
information and is therefore a valid source overcover. For a real-valued
path $Y$ and integers $a\le z$, put
\[
 \Var(Y;[a,z])=\sum_{e=a}^{z-1}|Y(e+1)-Y(e)|,
 \qquad \Var(Y)=\Var(Y;[0,T-1]).
\]
In particular, $\Var(x)\le T-1$.

\begin{definition}[Complete families and the two-class condition]
\label{def:complete-families}
\formal{SourceCovers/FamilyAccounting.lean}{SourceCovers.actualOwner,SourceCovers.ActualFamily,SourceCovers.HasSpanColoring,SourceCovers.mem_family_edges_iff,SourceCovers.constant_extension_isWalk}
A long hop $(p,u)$ has $u-p\ge2$. Its side is upper or lower according
as the work head stays strictly above or strictly below its endpoint cell
between $p$ and~$u$. Its owner is the first time strictly between the
endpoints at which its distance from that cell is largest. A
\emph{complete actual family} consists of all long hops of one physical
head, on one side, with one owner. For a family $G$, let $o_G$ be that
owner, let $(p_G,U_G)$ be its outermost hop, let
$I_G=(p_G,U_G]$ be its full half-open span, and let $h_G$ be its hop count.
Families are formed from all actual long hops before any selection or
payment. They are never redefined by deleting inconvenient hops.

The full family system has an \emph{$N$-class span partition} if it
partitions into $N$ classes whose spans are pairwise disjoint within each
class. A selected core is \emph{Q2} if its complete spans admit such a
partition into two classes. Empty classes are permitted. The full-system
constant $N$ need not be two.
\end{definition}

These are the first-maximum families of
\cite[Sections~2.2 and~3.2]{EighthPaper2026}. In that paper a bounded
owner-load supplies a bounded number of disjoint-span classes. Here the
existence of a fixed $N$-class partition is an explicit hypothesis; we
do not infer it from the number of physical tapes alone.

\begin{definition}[Admissible return geometry]\label{def:admissible-geometry}
An admissible geometry with parameters $d,b,c,\beta,N,q,a_0$ consists
of a fixed complete geometry and nonempty correctness set $\mathcal C$
as above, with $T\le cL$, whole-step capacity at most $\beta$, an
$N$-class partition of its full actual family system, an original Q2
core $\mathcal K_0$, and a disjoint set $\mathcal E$ of at most $q$
other complete families. Every family outside
$\mathcal K_0\cup\mathcal E$ has at most $a_0$ hops.
The local computation is deterministic; no linearity or product
assumption is imposed on $\mathcal C$. The word ``admissible'' below
includes all these premises. In a sequence, $c,\beta,N,q$ are fixed;
the geometry, core and exceptional families may vary with dimension.
\end{definition}

\begin{lemma}[Family accounting]\label{lem:family-accounting}
\formal{SourceCovers/FamilyAccounting.lean}{SourceCovers.family_owner_inside,SourceCovers.family_hops_nested,SourceCovers.family_span_hop,SourceCovers.family_size_le_owner_arms,SourceCovers.two_mul_sum_family_size_le,SourceCovers.sum_family_variation_eq,SourceCovers.q2_accounting,SourceCovers.mem_boundaryLowers,SourceCovers.boundary_lowers_le_crossings,SourceCovers.boundary_lowers_le_family_sizes,SourceCovers.boundary_lowers_le}
The hops in a family are strictly nested. If its ends are
$u_1<\cdots<u_{h_G}=U_G$, their starts satisfy
$p_1>\cdots>p_{h_G}=p_G$, and every hop contains $o_G$ strictly.
The outer span is itself a hop, and
\[
 h_G\le\min(o_G-p_G,U_G-o_G)\le\frac{U_G-p_G}{2}.
\]
For any set $\mathcal K$ of complete families, let
$m_{\mathcal K}(e)=\card{\{G\in\mathcal K:p_G\le e<U_G\}}$
be its overlap at the elementary boundary $e\mid e+1$. For every real path $Y$,
\begin{align}\label{eq:family-overlap-profile}
 2\sum_{G\in\mathcal K}h_G
 &\le\sum_{e=0}^{T-2}m_{\mathcal K}(e),\notag\\
 \sum_{G\in\mathcal K}\Var(Y;[p_G,U_G])
 &=\sum_{e=0}^{T-2}m_{\mathcal K}(e)|Y(e+1)-Y(e)|.
\end{align}
In particular, a Q2 core has $m_{\mathcal K}(e)\le2$, and
\begin{equation}\label{eq:q2-mass-variation}
 \sum_{G\in\mathcal K}h_G\le T-1,
 \qquad
 \sum_{G\in\mathcal K}\Var(x;[p_G,U_G])
       \le2\Var(x)\le2(T-1).
\end{equation}
For a set $\mathcal U$ of unselected families, let $Z_{\mathcal U}(e)$
be the distinct lower endpoints of its work edges and the state edge
crossing $e\mid e+1$, and let
$c_G(e)=\card{\{(p,u)\in G:p\le e<u\}}$. At every boundary inside the
computation,
\begin{equation}\label{eq:local-boundary-bill}
 |Z_{\mathcal U}(e)|
 \le1+\sum_{\substack{G\in\mathcal U\\ p_G\le e<U_G}}c_G(e)
 \le1+\sum_{\substack{G\in\mathcal U\\ p_G\le e<U_G}}h_G.
\end{equation}
Thus, if every unselected family has at most $a$ hops, this count is at
most $1+m_{\mathcal U}(e)a\le1+Na$. The same conclusion holds under
the local bound $m_{\mathcal U}(e)\le N$, without supplying a span partition.
\end{lemma}

\begin{outline}
The formalization imports the eighth paper's actual-hop, owner and
family definitions and their verified nesting, outer-span and height
bounds. Uniqueness of previous visits makes nesting strict. After
reflection for a lower family, its hop count is at most the height
of the common owner above the outer endpoint cell. The unit-step
distance bound limits that height separately by $o_G-p_G$ and
$U_G-o_G$, giving the owner-sensitive estimate.

Index an elementary edge by its lower time $e$. The full span
$(p_G,U_G]$ contains its upper time $e+1$ exactly when
$p_G\le e<U_G$. Double counting these incidences gives
\cref{eq:family-overlap-profile}: use unit edge weights for span
lengths, and $|Y(e+1)-Y(e)|$ for variation. The shared counting
identity works for weights in any additive commutative monoid.
Families active at one boundary have distinct colors in any
disjoint-span coloring. A Q2 core therefore has overlap at most two;
the input's unit-step bound gives \cref{eq:q2-mass-variation}.

For the boundary count, the Lean set contains the state lower
endpoint and the union of all crossing-hop lowers in the active
families. Its membership theorem proves that this restriction to
active families omits no crossing hop. The finite-union and image
cardinality bounds give \cref{eq:local-boundary-bill}. Applying the
family-size bound $a$ and overlap bound $N$ gives $1+Na$.
This last counting argument needs neither unit steps nor nesting
once the local overlap and size bounds are supplied.
\end{outline}

The source-bank and grid-payment arguments below use the uniform
specializations. The exact overlap and crossing-hop profiles retain
information that can reduce a geometry's payment estimate; no stronger
uniform exclusion is inferred from them here.

\subsection{Clipped shadows and exact core orbits}

\begin{definition}[Paid clipping and selected orbits]\label{def:clipped-shadow}
\formal{SourceCovers/ClippedShadows.lean}{SourceCovers.ActualFamily.shadow,SourceCovers.ActualFamily.shadow_at_hop,SourceCovers.ActualFamily.shadow_exact,SourceCovers.shadowOrbit}
Pay the outer lower endpoint $p_G$ of each selected family. Its
\emph{clipped shadow} $s_G$ is the identity off $(o_G,U_G]$. On that
interval it initially holds $o_G$; at each hop end $u_i$ it takes the
value $p_i$ and holds it until the next end, with every held value
replaced by its maximum with $p_G+1$. For a family set $\mathcal S$,
let $\mathcal O_{\mathcal S}(t)$ be the set of values obtained from $t$
by all finite compositions of its clipped shadows, including the empty
composition.
\end{definition}

The shadow is antitone on its active interval, and every nonidentity
action satisfies
\begin{equation}\label{eq:strict-shadow}
 p_G<s_G(t)\le o_G<t\le U_G.
\end{equation}
Clipping changes an actual hop value only at the outermost end $U_G$,
whose lower endpoint has been paid. Hence every selected actual hop
with an unpaid lower endpoint is represented exactly. We do not gate
these maps at other paid vertices. Their orbits may retain extra paths
through paid vertices, which is safe for a source overcover and preserves
closure under each selected shadow.

For a path $Y$ in a metric space, use the same variation notation with
$\operatorname{dist}(Y(e),Y(e+1))$ in place of the absolute difference.
This agrees with the earlier definition for real paths.

\begin{lemma}[Variation of a shadow or successful path]
\label{lem:shadow-variation}
\formal{SourceCovers/ClippedShadows.lean}{SourceCovers.jointShadow_eq_on,SourceCovers.jointShadow_variation,SourceCovers.shadow_success_cons_iff,SourceCovers.shadow_word_structure,SourceCovers.shadow_word_variation,SourceCovers.variation_comp_le_crossings,SourceCovers.exists_sharp_patch}
The joint map of clipped shadows with pairwise disjoint full spans has
\[
 \Var(Y\circ s)\le3\Var(Y)
\]
for every metric-valued path $Y$. For any fixed word of clipped shadows,
the set of starting times at which every action is nonidentity is an
integer interval, possibly empty. Its endpoint map is monotone in one
direction on that interval and no greater than its argument. Extending
that endpoint map by the identity outside the success interval also
gives variation at most $3\Var(Y)$ for every $Y$. The word may be empty
and may repeat families.

More generally, for a time map $f:\{0,\ldots,n\}\to\{0,\ldots,N\}$,
let $c_f(z)$ count the indices $e<n$ for which
$\min(f(e),f(e+1))\le z<\max(f(e),f(e+1))$. Then
\begin{equation}\label{eq:crossing-variation}
 \Var(Y\circ f;[0,n])
 \le\sum_{z=0}^{N-1}c_f(z)\operatorname{dist}(Y(z),Y(z+1)).
\end{equation}
The factor three is attained in the general class of regressive interval
patches, even by one constant active map.
\end{lemma}

\begin{outline}
The concrete held value is the minimum, over all hops in a complete
family, of the hop's start if it has ended and the owner otherwise.
It is antitone in time. Strict nesting identifies this minimum with
the actual start at each hop end. Taking its maximum with $p_G+1$
gives exactly the clipped shadow: every hop whose lower endpoint is
unpaid is preserved. Its active interval is exactly its nonidentity set.

For the reusable variation argument, sum the triangle inequality along
the elementary edges crossed by each jump, then interchange the finite
sums. The shared incidence identity gives \cref{eq:crossing-variation}.
A map monotone in either direction on one interval and equal to the
identity elsewhere has at most one downward crossing of each threshold.
Indeed, two downward crossings would force the intervening points into
the active interval. An antitone map could not rise between them; for a
monotone map, both downward crossings would have to occur at the two
joins, contradicting the order of the exterior identity values.
Telescoping the threshold's zero--one indicator bounds the total number
of crossings by twice the downward count plus one, hence by three.
This argument needs only preservation of the finite time range, not
regression of the patch.

For a disjoint packing, a threshold outside a family's full span is
unaffected: both the argument and its image stay on the same side.
At most one full span straddles a given threshold, so the packed map
has the same threshold indicator as that one shadow, or as the identity.
The three-crossing bound therefore holds independently of the family count.

Induct on a word. Its success domain is the first action's interval
intersected with the preimage of the tail's success domain. Mathlib's
order-convex preimage lemmas make this an interval. Composing the
monotone restrictions preserves one of the two orientations, and every
action weakly decreases time even when inactive. The successful endpoint
map therefore has the required interval, monotonicity and range properties.
Apply the single-patch bound to its identity extension. Empty success
domains and the empty word are included.

For sharpness of the general patch bound, take $f(t)=1$ on $(2,4]$ and
$f(t)=t$ elsewhere, on times $0,\ldots,5$. With $Y(t)=0$ for $t<2$ and
$Y(t)=1$ otherwise, $\Var(Y)=1$ and $\Var(Y\circ f)=3$. Lean
checks this finite witness directly. No sharpness claim for a particular
machine's source bank is inferred from this example.
\end{outline}

This lemma supplies the any-argument-path bound used repeatedly in the
exact source bank, without a factor exponential in successful word length.
The crossing profile can improve a particular map's estimate; the later
uniform constants still use the factor-three specialization. The packed
bank is formalized next; recovery through the residual dependencies remains
the subsequent task.

\begin{lemma}[The Q2 overlap forest and protected phases]
\label{lem:protected-phase}
\formal{SourceCovers/ProtectedPhases.lean}{SourceCovers.q2_caterpillar,SourceCovers.protected_survives,SourceCovers.admissiblePhases_card,SourceCovers.phase_payment_sum,SourceCovers.phase_average_exact,SourceCovers.protected_phase_average,SourceCovers.nonleaf_mass_le_total,SourceCovers.phase_path_length,SourceCovers.phase_success_length,SourceCovers.phase_shadowOrbit_exact,SourceCovers.phase_orbit_chains,SourceCovers.phase_orbit_card}
Join core families whose half-open spans intersect. For a Q2 core this
graph is a caterpillar forest: the vertices with at least two neighbors,
called the \emph{original nonleaves}, form a path within each component
where they occur. Let $\mathcal N$ be their set. Rank these paths from
left to right, starting at one separately in each component. All these
choices refer to the core before deletion.

Fix $M\ge2$ and a protected set $\mathcal P$ of at most $R<M$ core
families. Let $\mathcal A$ be the residues modulo $M$ not used by the
ranks of $\mathcal P\cap\mathcal N$, and let $F$ be the number of
distinct residues they use. For $\rho\in\mathcal A$, delete only the
original nonleaves of rank congruent to $\rho$, paying every actual
lower endpoint of their hops; call this union $\Gamma_{\rm rank}(\rho)$.
Every protected family survives. Put
$\tau_{\mathcal N}=\sum_{G\in\mathcal N\setminus\mathcal P}h_G$.
Then
\begin{equation}\label{eq:rank-phase-average}
 \card{\mathcal A}=M-F\ge M-R,
 \qquad
 \frac1{\card{\mathcal A}}
 \sum_{\rho\in\mathcal A}\card{\Gamma_{\rm rank}(\rho)}
 \le\frac{\tau_{\mathcal N}}{\card{\mathcal A}}
 \le\frac{\tau_{\mathcal N}}{M-R}.
\end{equation}
More precisely, the sum of paid-union sizes is at most
$\sum_{G\in\mathcal N\setminus\mathcal P}\card{L_G}$, where
$L_G$ is the set of actual hop lowers of $G$. The average bound with
denominator $M-F$ only requires $F<M$, regardless of $\card{\mathcal P}$.

In every phase, each surviving simple overlap path has at most $M+1$
vertices; in particular each successful core word has at most $M+1$
actions. The complete core orbit is the union of at most two deterministic
chains from the root and has at most $\ell=2M+3$ time values. Original
leaves and components of one or two vertices need not be deleted.
\end{lemma}

\begin{outline}
The overlap graph has no triangle: three pairwise intersecting intervals
have a common point, but one Q2 class would contain two of them. It has
no longer cycle either. On a cycle choose a vertex with largest left
endpoint. Its two cycle neighbors both extend strictly past that
endpoint, hence overlap each other, giving a triangle. Thus the graph
is a forest. Lean reuses Mathlib's cycle rotation and path operations.

An interval contained in another can have no further neighbor without
creating a triangle. Consequently nonleaf intervals have distinct left
endpoints and strictly increasing right endpoints in that order. An
overlap edge between nonleaves cannot skip a third nonleaf in the order:
that third interval would meet both endpoints and create a triangle.
Lean defines the rank as one plus the number of earlier nonleaves in
the same connected component, and proves that adjacent ranks differ
by one and distinct nonleaves have distinct ranks. Every internal vertex
of a simple path is a nonleaf, so removing original leaves preserves
connectivity between nonleaves. These facts give the connected nonleaf
path and its maximum degree two.

A finite unit-step integer chain avoiding one residue modulo $M$ visits
at most $M-1$ distinct integers. Lean proves this using the first forbidden
integer at or above the chain's minimum; a unit step cannot cross that
integer without visiting it. Apply this to the internal vertices of a
surviving simple path. Their ranks are distinct, so including the two
endpoints gives at most $M+1$ vertices. This argument needs no deletion
of leaves and also covers components of one or two vertices.

By \eqref{eq:strict-shadow}, consecutive successful actions have
overlapping spans and strictly decreasing owners. Their names form a
simple path, to which the preceding bound applies.
After an action from one Q2 class, its image lies in the same family's
span at or before its owner. That family is now inactive and no other
family of its class can act there. The only possible next action is
therefore the unique active family of the other class. At the root
there are at most two choices of first action; thereafter continuation
is deterministic. Including all intermediate stopping points gives
at most $1+2(M+1)=2M+3$ values.

The shared finite-orbit proof erases inactive actions from an arbitrary
word without changing its endpoint. It identifies the full word orbit
exactly with the values reachable in at most $M+1$ nonidentity steps.
This includes the empty word, intermediate stopping points and any
collisions between the two chains.

Only protected nonleaves forbid residues. Each unprotected nonleaf is
deleted in at most one admissible phase. Mathlib's finite-union bound
charges its actual lower support, whose size is at most its hop count.
Interchanging the finite sums gives the payment inequality, and counting
distinct forbidden residues gives $\card{\mathcal A}=M-F$.
\end{outline}

This supplies the finite core-orbit bound needed for the exact source
bank. Compared with deleting leaves together with their neighbors, it
retains the same orbit bound while paying only for original nonleaves.
The earlier estimate using
$\tau=\sum_{G\in\mathcal K\setminus\mathcal P}h_G$ still follows from
$\tau_{\mathcal N}\le\tau$. The next theorem assembles these bounds into a uniform packed bank.

\begin{theorem}[Theorem S: an exact source bank]\label{thm:exact-source-bank}
\formal{SourceCovers/ExactSourceBank.lean}{SourceCovers.exact_source_bank,SourceCovers.exactSourceBank_length,SourceCovers.exactSourceBank_bounded,SourceCovers.exactSourceBank_exact,SourceCovers.exactSourceBank_cover,SourceCovers.exactSourceBank_raw_root,SourceCovers.exactSourceBank_raw_values_card,SourceCovers.exactSourceBank_raw_selected_edge,SourceCovers.exactSourceBank_raw_variation}
Select a Q2 core of at most $J$ complete families, at most $q$ further
complete exceptional families, and a protected subcore of size at most
$R$. Let $M\ge\max(2,R+1)$ and use any admissible phase of
\cref{lem:protected-phase}. Pay the rank lowers and the outer lower
endpoints of all selected core and exceptional families. Put
\begin{equation}\label{eq:bank-constants}
 \begin{split}
  \ell&=2M+3,\qquad m_0=4+2M(M-1),\qquad C_0=16+24M(M-1),\\
  m_q&=m_0(1+m_0)^q,\qquad C_q=C_0(1+3C_0)^q.
 \end{split}
\end{equation}
There is a list $\mathcal B$ of at most $m_q$ geometry-fixed total time
maps such that, simultaneously for all time vertices $t$,
\begin{enumerate}
\item $\{a(t):a\in\mathcal B\}$ is exactly the complete orbit of the
      retained core and all selected exceptional shadows;
\item for every real-valued path $Y$,
      $\sum_{a\in\mathcal B}\Var(Y\circ a)\le C_q\Var(Y)$;
\item the time-value set is a union of at most $(\ell+1)^q$ complete
      core orbits, each a union of the two chains in
      \cref{lem:protected-phase}.
\end{enumerate}
Consequently the raw bank $Y_a=x\circ a$ contains the root read,
has total variation at most $C_q(T-1)$, and is set-closed across every
selected actual edge whose lower endpoint is unpaid. Its pointwise
cardinality is at most $\ell(\ell+1)^q$. For fixed $M$ and $q$, the
constants do not depend on $J$ or on the sizes or scales of the original
families.
\end{theorem}

\begin{outline}
Keep the original nonleaves as spine vertices. In each two-vertex
component, add its color-zero vertex to the spine; leave isolated vertices
in the leaf pack. The remaining leaves have pairwise disjoint spans,
so their joint shadow $s$ has variation constant three.
Every internal vertex of a successful word is an original nonleaf.
Thus packing leaves occur only at the two ends of such a word.

Label a surviving nonleaf by its original rank modulo $M$, excluding
the deleted residue. These $H=M-1$ labels are injective within each
surviving spine component: along a path, ranks make unit steps and avoid
the deleted residue, so they stay in an interval shorter than one period.
Equal residues therefore mean equal ranks and equal vertices. The extra
spine vertices from two-vertex components are isolated in the spine and
may use any of the $H$ labels. This uses $M\ge2$.

For each unordered pair of labels, with repetition allowed, pack one map
across all surviving spine components. In a component, choose the
nonempty successful word with those endpoint labels, if one exists.
The forest has a unique simple path between two vertices, and strictly
decreasing owners fix its orientation. Hence that word is unique.
Use its endpoint map on its success interval and identity elsewhere.
If no word exists, use identity throughout that component.

The union of the spans in a connected component is its interval hull,
and distinct component hulls are disjoint. Each chosen word remains in
its component hull. The shared disjoint-patch theorem therefore bounds
each packed map's variation by three, including both joins to identity.
Add an identity slot. Erasing inactive actions from any spine word shows
that these slots give exactly its complete orbit, simultaneously at every
root. Their number and summed variation coefficient are
\[
 m_S=1+\binom M2,
 \qquad C_S=1+3\binom M2.
\]
All choices precede the root and the argument path; a slot's representing
word is allowed to vary between components.

For each spine slot $a$, use
\[
 a,\qquad a\circ s,\qquad s\circ a,\qquad s\circ a\circ s.
\]
The two endpoint leaves and the spine middle of every successful core
word are represented by these four choices. Conversely each value follows
a core word at that argument. This proves exactness, rather than merely
coverage. The any-path variation bound composes: the four contributions
are at most $C_S,3C_S,3C_S,9C_S$. Thus there are $4m_S=m_0$
slots with total coefficient $16C_S=C_0$. The pointwise bound remains
$\ell$, by \cref{lem:protected-phase}.

Sort the exceptional families by decreasing owner, retaining a fixed tie
order. Successful words have strictly decreasing owners, so each such word
uses an ordered exceptional subsequence, with core words between its
members. Let $B$ be the core bank and $D$ the bank for the remaining
exception list. On inserting its first shadow $f$, form
\[
 D\ \mathbin{+\!+}\ \{b\circ f\circ a:b\in D,\ a\in B\},
\]
where $+\!+$ denotes list concatenation. Exactness follows by splitting a
successful word at $f$, or omitting $f$ if it is absent. Every constructed
value follows a selected word, even though a packed core slot need not
represent one fixed word globally. Slot count and summed variation obey
\[
 m_{r+1}=(1+m_0)m_r,
 \qquad C_{r+1}=(1+3C_0)C_r.
\]
These recurrences give \eqref{eq:bank-constants} without separately
formalizing the binomial sums over subsequences.

For the complete-core-orbit cover, retain the starting vertices of each
final core continuation. Each inserted exception creates at most $\ell$
new choices for each old continuation pattern, in addition to the pattern
omitting it. Induction gives at most $(\ell+1)^q$ terminal core roots.
The equality of the covered set with the enlarged orbit proves its
pointwise cardinality bound.

The empty word gives the root read. Prefixing a selected shadow gives
orbit set-closure. For an original core family removed by the phase,
every actual lower endpoint lies in the rank payment; it therefore has
no unpaid actual edge. For every remaining selected family, strict
clipping represents the actual hop exactly unless its outer lower was
paid. This proves closure across all originally selected unpaid edges,
and taking images under $x$ preserves it. Finally, the unit-step bound
$\Var(x)\le T-1$ gives the raw variation estimate. Other paid vertices
do not truncate the orbit.
\end{outline}

The formal variation bound holds for every pseudometric-valued path and
every finite horizon, with the same constants. The geometric bank needs
no unit-step assumption on the work walks; that assumption is used when
identifying shadows with actual hops. It also needs no numerical bound
on the number of core families. In the generic exceptional-list argument,
repeated entries and overlap with the core are allowed; the displayed
bounds count list entries.


Theorem~S supplies a representation of complete selected dependence
whose variation cost, for fixed $M$ and $q$, is independent of the number
of selected core families.
The next lemma pays for the remaining actual dependencies. Without this
step, Theorem~S alone is not a transposition certificate.

\subsection{Residual cuts and complete functional recovery}

For an integer width $w\ge1$ and shift
$\theta\in\{0,\ldots,w-1\}$, define the linear, endpoint-clipped block
\begin{equation}\label{eq:raw-block}
 B_\theta(y)=
 \{z\in\Omega:\lfloor(z+\theta)/w\rfloor
                   =\lfloor(y+\theta)/w\rfloor\}.
\end{equation}
Thus $|z-y|\le w-1$ within a block, including a partial block at either
end of the input.

\begin{lemma}[Paid residual-source recovery]\label{lem:residual-source-recovery}
\notformalized
Let $\Gamma_0$ be an initial paid set, and let $Y_1,\ldots,Y_m$ be
geometry-fixed maps from the timeline to $\Omega$. Their value set
$A(t)$ contains $x(t)$. Partition all dependency edges into selected
and residual edges. Suppose $A(u)\subseteq A(v)$ for every selected
edge $v\to u$ with $u\notin\Gamma_0$. For an unpaid lower endpoint
$u$, let $m_R(u)$ be its latest residual successor, or $u$ if none.
Define
\[
 V=\sum_i\Var(Y_i),\qquad
 \mathcal R=\sum_{u\notin\Gamma_0}\sum_i
   \operatorname{osc}(Y_i;[u,m_R(u)]),
\]
where oscillation is maximum minus minimum. If at most $H$ distinct
unpaid residual lowers are live across every elementary time boundary,
then, for every $w\ge1$, there are actual residual-lower sets
$\Gamma_\theta$ satisfying
\begin{equation}\label{eq:residual-average}
 \mathbb E_\theta\card{\Gamma_\theta}
       \le\frac{\mathcal R}{w}\le\frac{HV}{w}.
\end{equation}
For each shift, after paying
$\Gamma=\Gamma_0\cup\Gamma_\theta$, every output at time $t_r$ is
determined on $\mathcal C$ by the same transcript $G_\Gamma$ and the
bits in
\begin{equation}\label{eq:complete-bank-recovery}
 \bigcup_{i=1}^m B_\theta(Y_i(t_r)).
\end{equation}
No bound on full path depth, linearity or independence on a genie fibre
is needed.
\end{lemma}

\begin{informalproof}
For each unpaid residual lower $u$, pay $u$ if, for some $i$, the label
$\lfloor(Y_i(t)+\theta)/w\rfloor$ is not constant on
$[u,m_R(u)]$. This is an actual union of lower endpoints, so a lower
used by several dependencies is paid once. If the integer range of one
slot on an interval has oscillation $z$, a uniformly shifted grid cuts
that range with probability $\min(1,z/w)$. The union bound over slots
and lowers gives the first inequality of \eqref{eq:residual-average}.
Oscillation is at most variation on the interval. Summing that variation
over lower lifetimes charges an elementary step at most $H$ times,
giving $\mathcal R\le HV$.

A surviving residual edge preserves every slot label. A surviving
selected edge preserves inclusion of the sets of labels by the
hypothesis. Along any surviving backward path, induct from its last
vertex to its root: the label of its direct input read belongs to the
root's label set, because $x(t)\in A(t)$ at the last vertex. Therefore
every full surviving raw source belongs to
\eqref{eq:complete-bank-recovery}. This treats paths of arbitrary
length and all alternations of state and physical work edges.

Finally recursively evaluate the deterministic step functions. A
parent in $\Gamma$ is supplied by $G_\Gamma$; every other parent is
evaluated from its own parents and direct read. Time strictly decreases,
so this recursion terminates. Initial values and reads outside the input
are constants. All unfixed raw reads belong to the displayed union,
which therefore determines the output. The same reasoning starts at
a revealed output root and retains its direct read and all unpaid
parents; revealing the root was never used to delete them.
\end{informalproof}

For Theorem~S, take the surviving selected family edges as selected,
retain all other state and physical work edges as residual, and disregard
only edges already blocked at their lower endpoints. If every family
outside the selected core and the selected exceptional set has at most
$a$ hops,
\cref{lem:family-accounting} gives $H\le1+Na$. In each admissible phase
there is a shift with
\begin{equation}\label{eq:native-genie-ledger}
 \card\Gamma\le
 J+q+\card{\Gamma_{\rm rank}(\rho)}
       +\frac{(1+Na)C_q(T-1)}w.
\end{equation}
In particular a phase attaining the average in
\eqref{eq:rank-phase-average} permits replacing its rank payment by
$\tau_{\mathcal N}/(M-R)$, or the older upper bound $\tau/(M-R)$.
More generally, one may average a complete signed
reward-minus-payment expression over the admissible phases and their
raw-grid shifts. A reward chosen in one phase or shift cannot be paired
with a payment from another. Deleted selected families are blocked by
their paid lowers; they are not silently reclassified as small residual
families.

\subsection{Cyclic freezing of actual transpose demands}

View $\Omega$ as the integer points on a circle of circumference $L$.
For a nonempty subset $S$, list the circular distances between consecutive
distinct points as $g_1,\ldots,g_v$; their sum is $L$. A singleton has
one gap of length $L$. Define its point-gap area by
\begin{equation}\label{eq:gap-area}
 \gaparea(S)=\frac1{2L^2}\sum_{h=1}^v g_h^2.
\end{equation}
The blocks in \eqref{eq:raw-block} remain linear blocks. The circle is
used only for an overcover and a common ordering of raw coordinates.

\begin{theorem}[Theorem F: cyclic freezing]\label{thm:cyclic-freezing}
\notformalized
Let $m\ge1$ and let $Y_i:\{0,\ldots,T-1\}\to\Omega$, $1\le i\le m$,
be geometry-fixed bank maps. Let $Z$ be a
deterministic helper determined by the input, with $H(Z)\le g$ for the
uniform law on $\mathcal C$. Suppose every output is determined on
$\mathcal C$ by $Z$ and the raw block union in
\eqref{eq:complete-bank-recovery}, at one width $w$ and shift $\theta$. Put
$V=\sum_i\Var(Y_i)$. For output column $j$, let
$s_j=t_{jW}$ be its first append time and set
$S_j=\{Y_i(s_j):1\le i\le m\}$ and $A_j=\gaparea(S_j)$. Then
\begin{equation}\label{eq:cyclic-freezing}
 \Delta\ge
 W\sum_{j=0}^{d-1}A_j
       -\frac{2V}{d}-2(w-1)-\frac{W+1}{2}
       -g.
\end{equation}
In particular the bound holds with $g=\beta\card\Gamma$ for an actual
whole-step genie. If $u$ additional raw bits at geometry-fixed coordinates
are revealed jointly with that genie, it holds with
$g=\beta\card\Gamma+u$.
The bound permits arbitrary nonlinear recovery and a correlated
correctness set. Its area term uses the actual transpose demand stripes,
not an assumption that demands and sources are independent.
\end{theorem}

\begin{informalproof}
Write $e_j=t_{(j+1)W-1}$ and
$\delta_j=\sum_i\Var(Y_i;[s_j,e_j])$. These column intervals have
disjoint elementary edges, so $\sum_j\delta_j\le V$. Every source
block used during column $j$ lies in the fixed circular overcover
\[
 T_j=\{z\in\Omega:\dist_{\rm circ}(z,S_j)\le r_j\},
 \qquad r_j=w-1+\delta_j.
\]
Indeed, its coordinate differs from the corresponding current anchor
by at most $w-1$, and that anchor differs from its frozen value by at
most $\delta_j$. This covers entire blocks, not just their representatives.

Choose a uniformly random circular cut and an independent uniformly
random orientation. They define one total order on all raw coordinates,
shared by every column. Select a demanded coordinate $z$ of column $j$
when $z\notin T_j$ and $z$ precedes every coordinate in $T_j$. Every
recovery edge between selected coordinates then points forward in this
same order, so their induced recovery relation is acyclic.

Consider a complementary circular run of $T_j$ having $a$ integer
points. For its $v$-th point, the two orientations give selection
probabilities $v/L$ and $(a+1-v)/L$. Their average is the constant
$(a+1)/(2L)$ throughout the run. The selection probability is zero
on $T_j$. The sum of these constants over complementary runs is at
most $1/2$: the run points together with one distinct covered endpoint
per run number at most $L$. If $T_j=\Omega$, all these sums are empty.

For a frozen point gap of length $g$, thickening by integer radius $r_j$
leaves $(g-1-2r_j)_+$ integer points. The mean selection probability
over all $L$ raw coordinates is consequently
\[
 A_{j,r_j}=\frac1{2L^2}
       \sum_{g}(g-1-2r_j)_+\bigl((g-1-2r_j)_++1\bigr).
\]
For integers $g\ge1,r\ge0$,
\[
 g^2-(g-1-2r)_+\bigl((g-1-2r)_++1\bigr)\le(4r+1)g.
\]
When $g\ge2r+1$, expansion proves the inequality; otherwise its
left side is $g^2\le(4r+1)g$. Summing and using $\sum g=L$ gives
\begin{equation}\label{eq:thickening-area}
 A_{j,r_j}\ge A_j-\frac{2r_j+1/2}{L}.
\end{equation}
There is no factor for the number of slots or distinct frozen points.

The actual demanded coordinates in column $j$ form the stripe
\[
 D_j=\{iW+jb+k:0\le i<d,\ 0\le k<b\}.
\]
For every circular integer interval $I$,
\begin{equation}\label{eq:stripe-discrepancy}
 \left|\card(D_j\cap I)-\frac{\card I}{d}\right|
       \le b(1-1/d)\le b.
\end{equation}
To see this, the indicator of $D_j$ has period $W$ and one occupied
block of length $b$ per period. Its centered prefix count
$F(n)=\card(D_j\cap[0,n))-n/d$ is periodic and has oscillation
$b(1-1/d)$: it rises with slope $1-1/d$ on the occupied block and
falls with slope $-1/d$ on its complement. An interval discrepancy
is a difference of two prefix values, including intervals across the
circular seam.

Apply \eqref{eq:stripe-discrepancy} to the complementary runs, weighted
by their constant selection probabilities. Their constants sum to at
most $1/2$, so the expected number of selected demands in column $j$
is at least
\[
 W A_{j,r_j}-b/2
 \ge W A_j-\frac{2r_j}{d}-\frac1{2d}-\frac b2.
\]
Summing over columns and using
$\sum_j r_j\le d(w-1)+V$ gives expected selected size at least
the first four terms on the right of \eqref{eq:cyclic-freezing}.

For any realized common order, let $A$ be its selected coordinates.
Given $Z$ and the raw complement $\Omega\setminus A$, recover
the selected coordinates in reverse order using the acyclic relation.
For $X$ uniform on $\mathcal C$, it follows that
\[
 \log_2\card{\mathcal C}=H(X)
       \le H(Z)+H(X_{\Omega\setminus A}\mid Z)
       \le g+L-\card A.
\]
Average this inequality over the common cut and orientation to obtain
\eqref{eq:cyclic-freezing}. This entropy argument does not infer a
transcript cardinality bound from an entropy bound, and requires no
product distribution on a helper fibre. The genie and joint-revelation
special cases follow from subadditivity of entropy and the stated bit
capacities.
\end{informalproof}

Theorem~F charges the same genie that supplies the complete bank.
Own-source demands are rejected in its acyclic selection, so no separate
own-block collision loss is added. If $M,q,c$ are fixed, Theorem~S gives
$V=O(L)$, and the variation and stripe errors are $o(L)$ as
$d\to\infty$. The substantive remaining question is whether the
point-gap area exceeds the rank and residual payments; the following
sections prove this under explicit source-cover hypotheses.
